\documentclass[11pt,a4paper]{article}

% ===== 中文支持 =====
\usepackage[UTF8]{ctex}

% ===== 页面设置 =====
\usepackage[margin=2.2cm]{geometry}

% ===== 数学 =====
\usepackage{amsmath,amssymb}

% ===== 列表 =====
\usepackage{enumitem}

% ===== 表格 =====
\usepackage{array,booktabs}

% ===== 超链接 =====
\usepackage[colorlinks=true,linkcolor=blue,urlcolor=blue]{hyperref}

\title{\bfseries P57复习参考题 2 解答}
\author{}
\date{}

\begin{document}
	
	\maketitle
	
	
	
	\section*{复习参考题 2}
	
	\subsection*{复习巩固}
	
	\noindent\textbf{1. 解析}
	
	设 A 型号帐篷有 $x$ 顶，则 B 型号帐篷有 $(x+5)$ 顶，则
	\[
	\begin{cases}
		x>0,\ x\in \mathbf{N}^*,\\
		x+5>0,\\
		4x<48,\\
		0\leq 5x-48<5,\\
		3(x+5)<48,\\
		4(x+5)>48,
	\end{cases}
	\]
	解得 $9\leq x<11$。
	
	\noindent\textbf{2. 答案}
	
	(1) $<$ \quad (2) $>$ \quad (3) $>$
	
	(1) 解析：$\dfrac{1}{a}-\dfrac{1}{b}=\dfrac{b-a}{ab}>0$。因为 $a>b>0$，所以 $a>0$，$b>0$，$a-b>0$，所以 $\dfrac{b-a}{ab}<0$，即 $\dfrac{1}{a}<\dfrac{1}{b}$。
	
	(2) 解法一：$\dfrac{a}{c-a}-\dfrac{b}{c-b}=\dfrac{ac-bc+ab-bc}{(c-a)(c-b)}=\dfrac{(a-b)c}{(c-a)(c-b)}$。因为 $c>a>b>0$，所以 $c-a>0$，$c-b>0$，$a-b>0$，所以 $\dfrac{(a-b)c}{(c-a)(c-b)}>0$，即 $\dfrac{a}{c-a}>\dfrac{b}{c-b}$。
	
	解法二：因为 $c>a>b>0$，所以 $c-b>c-a>0$，所以 $\dfrac{1}{c-a}>\dfrac{1}{c-b}>0$，又 $a>b>0$，所以 $\dfrac{a}{c-a}>\dfrac{b}{c-b}$。
	
	(3) 解析：$\dfrac{a}{b}-\dfrac{a+c}{b+c}=\dfrac{ab+ac-ab-bc}{b(b+c)}=\dfrac{(a-b)c}{b(b+c)}$。因为 $a>b>c>0$，所以 $a-b>0$，$b+c>0$，所以 $\dfrac{(a-b)c}{b(b+c)}>0$，即 $\dfrac{a}{b}>\dfrac{a+c}{b+c}$。
	
	\noindent\textbf{3. 解析}
	
	(1) 设弧长为 $l$，半径为 $R$，则 $\dfrac{1}{2}lR=S$，周长 $y=2R+l\geq 2\sqrt{2Rl}=2\sqrt{4S}=4\sqrt{S}$，当且仅当 $2R=l$，即 $R=\sqrt{S}$ 时，等号成立，此时周长最小，为 $4\sqrt{S}$。
	
	(2) 设弧长为 $l$，半径为 $R$，则 $2R+l=P$，面积 $S=\dfrac{1}{2}lR\leq \dfrac{1}{2}R\left(\dfrac{R+l}{2}\right)^2=\dfrac{P^2}{16}$，当且仅当 $R=\dfrac{1}{2}l$，即 $R=\dfrac{P}{4}$ 时，等号成立，此时面积最大，为 $\dfrac{P^2}{16}$。
	
	\noindent\textbf{4. 解析}
	
	(1) $\left\{x\;\middle|\;-2\leq x\leq \dfrac{7}{4}\right\}$
	
	(2) $\{x\mid 5\leq x\leq 9\}$
	
	(3) $\mathbf{R}$
	
	(4) $\left\{x\;\middle|\;x<-\dfrac{1}{2}\ \text{或}\ x>1\right\}$
	
	\subsection*{综合运用}
	
	\noindent\textbf{5. 解析}
	
	因为 $a,b>0$，所以 $ab-3=a+b\geq 2\sqrt{ab}$，所以 $ab-2\sqrt{ab}-3\geq 0$，所以 $\sqrt{ab}\leq -1$（舍去）或 $\sqrt{ab}\geq 3$，所以 $ab\geq 9$，当且仅当 $a=b=3$ 时等号成立，所以 $ab$ 的取值范围是 $ab\geq 9$。
	
	\noindent\textbf{6. 解析}
	
	由题意可得 $k\neq 0$，$2kx^2+kx-\dfrac{3}{8}<0$ 对一切实数 $x$ 都成立，所以 $k<0$，且对于方程 $2kx^2+kx-\dfrac{3}{8}=0$，$\Delta=k^2-4\cdot 2k\left(-\dfrac{3}{8}\right)<0$，解得 $-3<k<0$。综上所述，$k$ 的取值范围是 $-3<k<0$。
	
	\noindent\textbf{7. 解析}
	
	(1) 设窗户面积为 $x$ m$^2$，则地板面积为 $(220-x)$ m$^2$。由题意可得
	\[
	\begin{cases}
		x<220-x,\\
		\dfrac{x}{220-x}\geq 10\%,\\
		x>0,\\
		220-x>0,
	\end{cases}
	\]
	解得 $20\leq x<110$，所以窗户面积至少为 $20$ m$^2$。
	
	(2) 公寓的采光效果变好了。设公寓原来的窗户面积和地板面积分别为 $x,y$（$x>0,y>0,x<y$），增加的面积为 $m$ m$^2$（$m>0$），由题意可得 $\dfrac{x+m}{y+m}-\dfrac{x}{y}=\dfrac{m(y-x)}{(y+m)y}>0$，即增加 $m$ m$^2$ 之后，比值变大了，故采光效果变好了。
	
	\noindent\textbf{8. 解析}
	
	\begin{center}
		\begin{tabular}{lll}
			\toprule
			\textbf{相等关系} & \textbf{不等关系} & \textbf{判断正误} \\
			\midrule
			(1) 若 $x=y$，则 $x^3=y^3$ & (1) 若 $x>y$，则 $x^3>y^3$ & 正确 \\
			(2) 若 $x=y$，则 $|x|=|y|$ & (2) 若 $x>y$，则 $|x|>|y|$ & 错误 \\
			(3) 若 $x=y$，则 $x^2=y^2$ & (3) 若 $x>y$，则 $x^2>y^2$ & 错误 \\
			(4) 若 $x=y\neq 0$，则 $\dfrac{1}{x}=\dfrac{1}{y}$ & (4) 若 $x>y>0$，则 $\dfrac{1}{x}>\dfrac{1}{y}$ & 错误 \\
			\bottomrule
		\end{tabular}
	\end{center}
	
	\subsection*{拓广探索}
	
	\noindent\textbf{9. 解析}
	
	设 $AM=y$ m，则 $x^2+4xy=200$，从而 $y=\dfrac{200-x^2}{4x}$。于是
	\[
	S=4\cdot 200x^2+210\times 4xy+80\times 2y^2
	\]
	\[
	=4\cdot 200x^2+210\times 4x\cdot \frac{200-x^2}{4x}+80\times 2\left(\frac{200-x^2}{4x}\right)^2
	\]
	\[
	=38000+4000x^2+\frac{400000}{x^2}\geq 38000+2\sqrt{4000x^2\cdot \frac{400000}{x^2}}=118000,
	\]
	当且仅当 $4000x^2=\dfrac{400000}{x^2}$，即 $x=\sqrt{10}$（$x=-\sqrt{10}$ 舍去）时等号成立。由上述可知，当 $AD=\sqrt{10}$ m 时，休闲广场总造价 $S$ 取最小值，为 $118000$ 元。
	
	\noindent\textbf{10. 解析}
	
	按第一种策略购物，设第一次购物时价格为 $P_1$ 元，购 $n$ kg，第二次购物时价格为 $P_2$ 元，仍购 $n$ kg（$P_1>0,P_2>0,n>0$），两次购物的平均价格为 $\dfrac{P_1n+P_2n}{2n}=\dfrac{P_1+P_2}{2}$；若按第二种策略购物，第一次花 $m$ 元，能购 $\dfrac{m}{P_1}$ kg 物品，第二次仍花 $m$ 元，能购 $\dfrac{m}{P_2}$ kg 物品，两次购物的平均价格为 $\dfrac{2m}{\dfrac{m}{P_1}+\dfrac{m}{P_2}}=\dfrac{2P_1P_2}{P_1+P_2}$。
	
	比较两次购物的平均价格。
	
	解法一：
	\[
	\frac{P_1+P_2}{2}-\frac{2P_1P_2}{P_1+P_2}
	=\frac{(P_1+P_2)^2-4P_1P_2}{2(P_1+P_2)}
	=\frac{(P_1-P_2)^2}{2(P_1+P_2)}\geq 0,
	\]
	所以 $\dfrac{P_1+P_2}{2}\geq \dfrac{2P_1P_2}{P_1+P_2}$，当且仅当 $P_1=P_2$ 时等号成立。
	
	解法二：
	\[
	\frac{2P_1P_2}{P_1+P_2}\leq \frac{2P_1P_2}{2\sqrt{P_1P_2}}=\sqrt{P_1P_2}\leq \frac{P_1+P_2}{2},
	\]
	当且仅当 $P_1=P_2$ 时等号成立。
	
	所以用第二种购物策略比较经济。
	
	一般地，如果 $n$ 次购买同一种物品，用第二种策略购买比较经济。
	
\end{document}